% ============================================================
% Chapter 7: Results
%   Section: Pilot on X-Fact
%     Subsection: What Was Run
% Included from chapters/07-results/01-xfact-pilot/xfact-pilot.tex
% ============================================================

\subsection{What Was Run}
\label{subsec:pilot-setup}

Nothing was trained. The verification model is a zero-shot language
model reading evidence that the pipeline retrieves from the open web
at the moment of checking, so there are no weights to fit and no
X-Fact claim was used to tune any threshold. What the pilot measures
is the pipeline's own claim-checking stage (the one behind the
single-claim endpoint), run through the evaluation harness of
Section~\ref{sec:harness} over X-Fact rows and scored against their
labels.

Two sets were run on 13 October 2026:

\begin{itemize}
    \item \textbf{The pilot}: 64 claims sampled from the English and
    Spanish rows of X-Fact, eight per language and label (fewer where
    a label is missing), preferring development and test rows. Its 24
    English claims come from PolitiFact and are all training rows,
    because English has no evaluation split in X-Fact; its 40 Spanish
    claims come from Chequeado (31 training, 2 development, 7 test).
    The label distribution is flat by construction.
    \item \textbf{The Spanish test split}: all 165 Spanish rows of
    X-Fact's in-domain test set ($\alpha_1$), extracted in file order
    by a deterministic script. It is the only part of the benchmark
    the paper reports on that this work can reproduce, and it covers
    one of $\alpha_1$'s twelve languages.
\end{itemize}

\paragraph{Configuration.} \texttt{llama3.2:3b} (Q4\_K\_M) served by
Ollama on the development laptop (Intel i7-12700H, 16~GB RAM, an
NVIDIA RTX~3050 Laptop GPU with 4~GB, of which about 80\% of the model
fitted, the rest running on the CPU); the inference service (entities
and embeddings) on the CPU in its container; SearXNG with the engine
set of the deployment. Thresholds and scoring weights were the
defaults, and the internal corpus was a snapshot of the live store
(5 points, none of them about the claims). Each claim was checked
without its claimant, dates or fact-checker: the claim text alone, as
the single-claim endpoint receives it.

\paragraph{Health before the run.} The health endpoint answered for
all four dependencies (inference 49~ms, SearXNG 36~ms, the language
model 42~ms, Neo4j 1,060~ms). The source probe's search check returned
89 results in each language, from Bing, Yep, Wikipedia (English only)
and four academic indexes; Google and Brave were refusing SearXNG
(``access denied'', ``too many requests''). A three-claim smoke run
returned 15 to 17 web results per claim before the full run started.
